\section{From the internal complex structure to the projective geometry of Bloch}
\label{sec:geometria-proyectiva-bloch-rev2}
\index{Bloch sphere}
\index{complex structure!orthogonal realification}
\index{binary icosahedral group}
\index{Möbius!projective action}

The internal square root of minus the identity has two coordinated realizations. In
Part I, Paley--Witt incidence constructs
\(J_W^2=-I\) over \(\mathbb F_3^6\), organizes the ternary code, and realizes
the structure \(\mathbb F_9^3\); in an analytic realization, an orthogonal operator
\(J_E\) with \(J_E^2=-I\) turns a real space of even
dimension into a complex space. This second arrow adds topology, norm, and
completion: it preserves the quadratic relation but does not identify the scalar
fields.

The complete transition is fully characterized by
\begin{equation}
 \boxed{
 \begin{gathered}
 C_P\longrightarrow J_W^2=-I_6
 \longrightarrow \mathbb F_3^6\simeq\mathbb F_9^3,\\
 (E_4,J_E),\quad \dim_{\mathbb R}E_4=4,
 \longrightarrow
 \mathbb P(\mathcal H_J)\simeq\mathbb {CP}^{1}
 \longrightarrow S^2,\\
 \mathrm{SU}(2)\longrightarrow\mathrm{PSU}(2)\simeq\mathrm{SO}(3)
 \subset\mathrm{PSL}(2,\mathbb C)\cong\text{\rmfamily Möb}^{+}(S^2),\\
 2A_5=\pi^{-1}(A_5)\subset\mathrm{SU}(2).
 \end{gathered}}
 \label{eq:cadena-paley-bloch-rev2}
\end{equation}
The first line retains its proofs in the
\cref{chap:paley-codigo-witt,thm:f9-tres-desde-paley}; the two projective lines are constructed here.

\subsection{Phase and ray space realification}

Let \((E,g)\) a space euclideo real of dimension par and let
\(J_E\in\operatorname{End}(E)\) such that
\[
 J_E^2=-I,
 \qquad g(J_Eu,J_Ev)=g(u,v).
\]
Then \(J_E^{\mathsf T}=-J_E\) and
\[
 R_J(\theta):=e^{\theta J_E}
 =\cos\theta\,I+\sin\theta\,J_E
\]
is orthogonal, with
\(R_J(\theta)R_J(\phi)=R_J(\theta+\phi)\). The action
\((a+ib)u:=au+bJ_Eu\) realizes the complex chart without using the scalar
\(i\) as antecedent of the operator.

Now let \(E_4\subseteq E\) be a \(J_E\)-stable subspace of real dimension
four and let \(\mathcal H_J=(E_4,J_E)\) be the resulting Hilbert space, of
complex dimension two.
A unit vector
\[
 \psi=\begin{pmatrix}z_0\\z_1\end{pmatrix},
 \qquad |z_0|^2+|z_1|^2=1,
\]
determines to ray \([\psi]\in\mathbb {CP}^{1}\). Definamos
\begin{equation}
 r_1=2\operatorname{Re}(\overline z_0z_1),
 \qquad
 r_2=2\operatorname{Im}(\overline z_0z_1),
 \qquad
 r_3=|z_0|^2-|z_1|^2.
 \label{eq:coordenadas-bloch-rev2}
\end{equation}

\begin{theorem}[Projective Bloch Identification]
\label{thm:bloch-proyectivo-rev2}
The map
\[
 \mathcal B:\mathbb {CP}^{1}\longrightarrow S^2,
 \qquad [\psi]\longmapsto(r_1,r_2,r_3),
\]
is a smooth bijection. The rank one projector associated with the ray is
\begin{equation}
 \rho_\psi=|\psi\rangle\langle\psi|
 =\frac12\bigl(I_2+r_1\sigma_1+r_2\sigma_2+r_3\sigma_3\bigr),
 \label{eq:proyector-bloch-rev2}
\end{equation}
where \(\sigma_1,\sigma_2,\sigma_3\) are the Pauli matrices.
\end{theorem}

\begin{proof}
The global phase leaves the three coordinates invariant and
\[
 r_1^2+r_2^2+r_3^2
 =4|z_0|^2|z_1|^2+(|z_0|^2-|z_1|^2)^2=1.
\]
Conversely, each \(\boldsymbol r\in S^2\) determines the positive projector
\(\rho=(I+\boldsymbol r\cdot\boldsymbol\sigma)/2\), which satisfies
\(\rho^2=\rho\) and \(\operatorname{tr}\rho=1\); its image is a single ray.
The affine charts of \(\mathbb {CP}^{1}\) make the map and its inverse
\cite{Bloch1946} smooth.
\end{proof}

The Hopf sequence
\begin{equation}
 S^1\longrightarrow S^3\longrightarrow S^2
 \label{eq:fibrado-hopf-bloch-rev2}
\end{equation}
exhibits the loss published by projectivization: the normalized vector
belongs to \(S^3\), the projective observable publishes a point of \(S^2\), and
the fiber \(S^1\) preserves the common phase that the chart identifies
\cite{Hopf1931}. This is a precise realization of the HMT distinction between
enriched state and visible coordinate.

\subsection{Rotations, Möbius transformations and spinor lifting}

Every element of \(\mathrm{SU}(2)\) can be written
\[
 U=\begin{pmatrix}
 \alpha&\beta\\-\overline\beta&\overline\alpha
 \end{pmatrix},
 \qquad |\alpha|^2+|\beta|^2=1.
\]

\begin{theorem}[Spinorial action on the Bloch sphere]
\label{thm:su2-so3-bloch-rev2}
For each \(U\in\mathrm{SU}(2)\) there exists a unique rotation
\(R_U\in\mathrm{SO}(3)\) such that
\[
 U(\boldsymbol r\cdot\boldsymbol\sigma)U^\dagger
 =(R_U\boldsymbol r)\cdot\boldsymbol\sigma.
\]
The map \(U\mapsto R_U\) is a surjective homomorphism with kernel
\(\{\pm I_2\}\).
\end{theorem}

\begin{proof}
Conjugation preserves the three-dimensional real space of traceless
Hermitian matrices and the inner product
\(\langle A,B\rangle=\tfrac12\operatorname{tr}(AB)\), and therefore induces a
rotation. If it acts trivially, \(U\) commutes with the three Pauli matrices
and is scalar; the unit-determinant condition leaves \(U=\pm I_2\). The axis--angle
representation provides the two preimages of each rotation.
\end{proof}

In the affine chart \(z=z_0/z_1\) of \(\mathbb {CP}^{1}\), the same action is
the Möbius transformation
\[
 z\longmapsto
 \frac{\alpha z+\beta}{-\overline\beta z+\overline\alpha}.
\]
The Bloch rotation and the fractional transformation are two coordinates of
a single projective action: the first acts on Hermitian projectors and
the second on the Riemann sphere.

The rotational group of the icosahedron is \(A_5\subset\mathrm{SO}(3)\). Its preimage defines the binary icosahedral group
\[
 2A_5:=\pi^{-1}(A_5)\subset\mathrm{SU}(2),
\]
and produces the exact sequence
\begin{equation}
 1\longrightarrow\{\pm I_2\}
 \longrightarrow2A_5\longrightarrow A_5\longrightarrow1.
 \label{eq:extension-central-2a5-rev2}
\end{equation}
In particular, \(|2A_5|=120\). For a rotation of odd order \(n\), the
half-angle lift has order \(2n\), whereas its opposite has
order \(n\). Icosahedral rotations of orders three and five therefore admit
pairs of lifts of orders \((6,3)\) and \((10,5)\),
respectively. This specialization links the
projective geometry of two-level states with the McKay branch and the
exceptional symmetries, whose modules and lattices remain in their
reference constructions in Part I.

\begin{statusbox}[title={Provenance and scope of the projective chain}]
The Paley--Witt identity, structure
\(\mathbb F_3^6\simeq\mathbb F_9^3\), action
\(\mathrm{SU}(2)\to\mathrm{SO}(3)\), the Mobius realization, and extension
\(2A_5\) are results demonstrated in their domains. Chain
\eqref{eq:cadena-paley-bloch-rev2} is a unified functorial realization:
it explicitly declares the realization that changes scalars and does not convert
the finite root \(J_W\) into an analytic operator without that additional datum.
\end{statusbox}
